\documentclass[11pt]{article}
\usepackage[utf8]{inputenc}
\usepackage{amsmath,amssymb,amsthm}
\usepackage[margin=1in]{geometry}
\usepackage{hyperref}
\usepackage{xcolor}

\newcommand{\tideref}[2][]{\href{https://therisensea.org/tides/#2/}{\texttt{#2}}}

\title{Tide retrospective: the finite-order smooth recovery\\
(smooth-germ programme, stages C4--C5)}
\author{Tide loop (Claude Fable 5, with GPT-5.6 Sol deliberation)}
\date{2026-08-10}

\begin{document}
\maketitle

\section{Setting}

The closing tide of the smooth-germ programme, and with it the
germbij nondegenerate arc: the bridge from the polynomial recovery
machinery to smooth losses, which is where Theorem~3.1 actually
lives. The consult's staging held: with the local Taylor comparison
in place, C4 was composition plus one genuinely new ingredient ---
which turned out to be a Mathlib lookup
(\texttt{taylor\_isLittleO}, the Peano remainder).

\section{The thing formalised}

The headline\footnote{\texttt{smooth\_jet\_recovery}.}:
\begin{quote}
Two admissible $C^{R+2}$ losses whose normalized moments at
$x^{2+r}$ differ by $o(q^{2+2r})$ for every $r \le R$ have the same
Taylor data through degree $R+2$: the quadratic coefficient
$\lambda/2$ and every higher coefficient.
\end{quote}
Supporting cast, all in \texttt{Laplace/OneD/SmoothRecovery.lean}:
the Peano remainder in epsilon-radius form; the derived minimum facts
(the derivative vanishes at an admissible minimum, and
$\rho \le \lambda/2$, so nondegeneracy costs no hypothesis); the
Taylor polynomial rewritten in jet shape; the stabilizer extension
and its identity; the stabilized jet's joint
admissibility-and-profile certificate; the loss-versus-stabilized-jet
epsilon bound (Peano plus stabilizer absorption at radius
$\varepsilon/(2(d+1))$); and the truncated strong induction.

\section{The one structural discovery}

The recovery theorem of the weighted-jet programme quantifies over
\emph{every} rung of its coefficient vector. The stabilized jet is
deliberately longer than the data reaches --- $M - 2 > R$ rungs, the
stabilizer living above the Taylor degree precisely so it cannot
disturb the recovered coefficients --- so the full-length theorem is
inapplicable, and honestly so: the data says nothing about the
stabilizer coefficients. The fix is
\texttt{jet\_recovery\_stable\_partial}, the same strong induction
truncated at rung $R$. This is the kind of quantifier mismatch that
only surfaces at composition time, and it is exactly the ``full-jet
quantifier management'' the original scoping consult listed among
C's obligations.

\section{Proof notes}

The per-rung assembly follows the recorded plan: the jet-moment
difference at observable $x^{2+r}$, scaled by $q^{-r}$, decomposes
through the scaling bridge into three pieces --- two
loss-versus-polynomial comparisons (each a C3 limit times
$q^{D-2-r}$, vanishing or bounded), and the data hypothesis verbatim
in the middle. The pow-splitting rewrites
($q^{2+r+D-2} = q^{2+2r} \cdot q^{D-2-r}$, proven by \texttt{omega}
on the exponents) keep natural-number subtraction away from
\texttt{ring}. One tactical lesson worth the catalogue: the ascribed
type of a large composed \texttt{Tendsto} term cannot realistically
match the term-mode composite's exact shape (associativity and
negation placement); dropping the ascription and letting
\texttt{congr'} carry the function shape --- with the eventual
identity closed by \texttt{field\_simp} and \texttt{ring} --- is the
robust pattern.

\section{Deliberation and check}

Prior-deliberation path: C4--C5 were specified in the scoping consult
and the composition plan was recorded in the tide log before writing.
No new closed form entered; the analytic content was checked at C3
(statement-level) and across the weighted-jet programme's executed
checks.

\section{What this closes, and what remains}

With this tide the germbij nondegenerate programme is formalised end
to end: identifiability (Theorem~7.3, 1D and $\mathbb{R}^d$), the
anharmonic and weighted-jet recovery ladders, Proposition~4.1 with
its witness, the polynomial comparison core with unequal bases, and
now the smooth finite-order recovery. The note's remaining
commissions are the multivariate Theorem~3.1 (the Susceptibility
Primer's missing Stage~3, needing Wick machinery on $\mathbb{R}^d$)
and the full-jet quantifier packaging (C6), both needing fresh
scoping. The natural next act after the merge is the note-side
markers on Theorem~3.1 for \texttt{smooth\_jet\_recovery}.

\end{document}
